\documentclass[11pt,a4paper]{article}
\usepackage[utf8]{inputenc}
\usepackage[T1]{fontenc}
\usepackage{amsmath,amssymb,amsthm}
\usepackage{geometry}
\geometry{margin=2.5cm}
\theoremstyle{plain}
\newtheorem{theorem}{Theorem}[section]
\newtheorem{proposition}[theorem]{Proposition}
\newtheorem{lemma}[theorem]{Lemma}
\theoremstyle{definition}
\newtheorem{definition}[theorem]{Definition}
\title{Soluciones Tests 6-10: Álgebra}
\author{Mathematical AI Agent}
\date{\today}
\begin{document}
\maketitle

\section{Problemas}

\begin{theorem}[Desigualdad de Cauchy-Schwarz]
Sea $a, b, c > 0$. Entonces
\[
(a + b + c)\left(\frac{1}{a} + \frac{1}{b} + \frac{1}{c}\right) \geq 9.
\]
\end{theorem}

\begin{proof}
Aplicamos la desigualdad de Cauchy-Schwarz en su forma general. Para cualquier secuencia de pares reales $(x_i, y_i)$, se cumple
\[
\left(\sum_{i=1}^{n} x_i^2\right)\left(\sum_{i=1}^{n} y_i^2\right) \geq \left(\sum_{i=1}^{n} x_i y_i\right)^2.
\]
Tomamos $n = 3$, $x_1 = \sqrt{a}$, $x_2 = \sqrt{b}$, $x_3 = \sqrt{c}$, $y_1 = 1/\sqrt{a}$, $y_2 = 1/\sqrt{b}$, $y_3 = 1/\sqrt{c}$. Entonces:
\[
\sum_{i=1}^{3} x_i^2 = a + b + c, \qquad \sum_{i=1}^{3} y_i^2 = \frac{1}{a} + \frac{1}{b} + \frac{1}{c},
\]
y
\[
\sum_{i=1}^{3} x_i y_i = \sqrt{a}\cdot\frac{1}{\sqrt{a}} + \sqrt{b}\cdot\frac{1}{\sqrt{b}} + \sqrt{c}\cdot\frac{1}{\sqrt{c}} = 1 + 1 + 1 = 3.
\]
Sustituyendo en la desigualdad de Cauchy-Schwarz obtenemos
\[
(a + b + c)\left(\frac{1}{a} + \frac{1}{b} + \frac{1}{c}\right) \geq 3^2 = 9.
\]
La igualdad se alcanza si y solo si $x_i/y_i$ es constante para todo $i$, es decir, $\sqrt{a}/(1/\sqrt{a}) = \sqrt{b}/(1/\sqrt{b}) = \sqrt{c}/(1/\sqrt{c})$, lo que equivale a $a = b = c$.
\end{proof}

\begin{theorem}[Cota en el círculo unitario]
Sean $x, y \in \mathbb{R}$ tales que $x^2 + y^2 = 1$. Entonces
\[
|x + y| \leq \sqrt{2}.
\]
\end{theorem}

\begin{proof}
Consideramos la expresión $(x + y)^2$. Expandiendo obtenemos
\[
(x + y)^2 = x^2 + 2xy + y^2 = (x^2 + y^2) + 2xy = 1 + 2xy.
\]
Por la desigualdad de la media aritmética-geométrica, para cualesquiera $x, y \in \mathbb{R}$ se tiene
\[
2xy \leq x^2 + y^2 = 1,
\]
puesto que $(x - y)^2 \geq 0$ implica $x^2 - 2xy + y^2 \geq 0$, de donde $2xy \leq x^2 + y^2$. Por lo tanto
\[
(x + y)^2 = 1 + 2xy \leq 1 + 1 = 2.
\]
Tomando raíz cuadrada de ambos lados (ya que $|x+y| \geq 0$):
\[
|x + y| = \sqrt{(x + y)^2} \leq \sqrt{2}.
\]
La igualdad se alcanza cuando $x = y = \pm 1/\sqrt{2}$.
\end{proof}

\begin{theorem}[Identidad con suma nula]
Sean $x, y, z \in \mathbb{R}$ tales que $x + y + z = 0$. Entonces
\[
x^2 + y^2 + z^2 = -2(xy + xz + yz).
\]
\end{theorem}

\begin{proof}
Partimos de la hipótesis $x + y + z = 0$ y elevamos al cuadrado:
\[
(x + y + z)^2 = 0^2 = 0.
\]
Expandiendo el cuadrado del trinomio:
\[
x^2 + y^2 + z^2 + 2xy + 2xz + 2yz = 0.
\]
Agrupando los términos lineales:
\[
x^2 + y^2 + z^2 + 2(xy + xz + yz) = 0.
\]
Despejando $x^2 + y^2 + z^2$:
\[
x^2 + y^2 + z^2 = -2(xy + xz + yz).
\]
\end{proof}

\begin{theorem}[Igualdad de la traza]
Sean $A, B \in M_n(\mathbb{R})$ matrices cuadradas de dimensión $n$. Entonces
\[
\operatorname{tr}(AB) = \operatorname{tr}(BA).
\]
\end{theorem}

\begin{proof}
Sean $A = (a_{ij})$ y $B = (b_{ij})$ matrices de orden $n \times n$. Por definición, la traza de una matriz es la suma de sus elementos diagonales:
\[
\operatorname{tr}(AB) = \sum_{i=1}^{n} (AB)_{ii}.
\]
El elemento $(i,i)$ del producto $AB$ se calcula como
\[
(AB)_{ii} = \sum_{j=1}^{n} a_{ij} b_{ji}.
\]
Por lo tanto,
\[
\operatorname{tr}(AB) = \sum_{i=1}^{n} \sum_{j=1}^{n} a_{ij} b_{ji}.
\]
De manera análoga,
\[
\operatorname{tr}(BA) = \sum_{i=1}^{n} (BA)_{ii} = \sum_{i=1}^{n} \sum_{j=1}^{n} b_{ij} a_{ji}.
\]
Interambiando los roles de los índices $i$ y $j$ en la última expresión (ya que son índices mudos de suma):
\[
\operatorname{tr}(BA) = \sum_{j=1}^{n} \sum_{i=1}^{n} b_{ji} a_{ij} = \sum_{i=1}^{n} \sum_{j=1}^{n} a_{ij} b_{ji} = \operatorname{tr}(AB).
\]
\end{proof}

\begin{theorem}[Determinante de la inversa]
Sea $A \in M_n(\mathbb{R})$ una matriz invertible. Entonces
\[
\det(A^{-1}) = \frac{1}{\det(A)}.
\]
\end{theorem}

\begin{proof}
Sea $A$ una matriz invertible de orden $n$. Entonces existe $A^{-1}$ tal que
\[
AA^{-1} = A^{-1}A = I_n,
\]
donde $I_n$ es la matriz identidad de orden $n$. Por la multiplicatividad del determinante, para cualesquiera matrices cuadradas $A$ y $B$ de la misma dimensión se cumple
\[
\det(AB) = \det(A)\cdot\det(B).
\]
Aplicando esta propiedad al producto $AA^{-1}$:
\[
\det(AA^{-1}) = \det(A)\cdot\det(A^{-1}).
\]
Como $AA^{-1} = I_n$ y sabemos que $\det(I_n) = 1$ (puesto que la matriz identidad tiene todos sus elementos diagonales iguales a 1 y los demás nulos, su determinante es el producto de estos elementos diagonales), obtenemos
\[
\det(A)\cdot\det(A^{-1}) = 1.
\]
Como $A$ es invertible, $\det(A) \neq 0$, por lo que podemos despejar:
\[
\det(A^{-1}) = \frac{1}{\det(A)}.
\]
\end{proof}

\end{document}
